\chapter{Geometria algebraiczna a topologia: od cykli do twierdzeń Lefschetza}
\label{ch:ga10-topologia}

\section{Dwa sposoby patrzenia na tę samą przestrzeń}
\label{sec:ga10-punkt-wyjscia}

Równanie wielomianowe wyznacza przestrzeń algebraiczną, ale nad $\C$
można badać również jej zwykłą topologię. Dla gładkiej odmiany
rzutowej $X$ nad $\C$ analityfikacja $X^{\rm an}$ jest zwartą,
zorientowaną rozmaitością gładką wymiaru rzeczywistego $2n$,
gdzie $n=\dim_{\C}X$. Orientację wyznaczają zespolone układy
współrzędnych: baza $(v_1,Jv_1,\ldots,v_n,Jv_n)$ jest dodatnio
zorientowana. To właśnie orientacja umożliwia całkowanie form
najwyższego stopnia i definiowanie liczb przecięć.

Chcemy odpowiedzieć na trzy pytania: jak z równań wielomianowych
obliczać grupy kohomologii; dlaczego klasy podrozmaitości tworzą
szczególne klasy kohomologiczne; jakie topologie są niemożliwe
dla gładkich odmian rzutowych?

W całym rozdziale $H^k(M;R)$ oznacza kohomologię singularną
przestrzeni topologicznej $M$ ze współczynnikami w pierścieniu $R$.
Jeśli piszemy $H^k_{\rm dR}(M)$, chodzi o kohomologię kompleksu
gładkich form różniczkowych o rzeczywistych współczynnikach.
Wymiary $b_k(M)=\dim_{\Q}H^k(M;\Q)$ nazywamy \emph{liczbami Bettiego}.
Dla zwartej rozmaitości są one skończone. Zapis $H^q(X,\mathcal F)$
nadal oznacza \emph{kohomologię snopa}; nie utożsamiamy tych dwóch
pojęć bez użycia odpowiedniego twierdzenia porównawczego.

\begin{przyklad}[Krzywa algebraiczna jako powierzchnia topologiczna]
Gładka spójna krzywa rzutowa $C$ nad $\C$ jest zwartą, spójną,
zorientowaną powierzchnią rzeczywistą. Jej rodzaj $g$ oznacza liczbę
uchwytów, a jej liczby Bettiego wynoszą $b_0=b_2=1$ i $b_1=2g$.
Dlatego $\chi_{\rm top}(C)=2-2g$; wyjaśnimy poniżej, jak ta
równość wynika zarówno z topologii, jak i z rachunku dywizorów.
\end{przyklad}

\section{Kohomologia de Rhama, dualność Poincarégo i cykle algebraiczne}
\label{sec:ga10-de-rham}

\begin{definicja}[Kohomologia de Rhama]
Niech $M$ będzie gładką rozmaitością rzeczywistą.
Dla $k\ge0$ przez $\Omega^k(M)$ oznaczamy przestrzeń
rzeczywistych gładkich $k$-form, a przez $\mathrm d\colon\Omega^k(M)
\to\Omega^{k+1}(M)$ różniczkę zewnętrzną. Ponieważ $\mathrm d^2=0$,
definiujemy
\[
 H^k_{\rm dR}(M)=\frac{\ker(\mathrm d\colon\Omega^k(M)\to
 \Omega^{k+1}(M))}{\operatorname{im}(\mathrm d\colon
 \Omega^{k-1}(M)\to\Omega^k(M))},
\]
z umową $\Omega^{-1}(M)=0$. Formy z licznika nazywamy
\emph{zamkniętymi}, a formy z mianownika \emph{dokładnymi}.
\end{definicja}

\begin{uwaga}[Dlaczego iloraz ma sens geometryczny?]
Formy dokładne mają zerową całkę po każdym cyklu bez brzegu:
na mocy twierdzenia Stokesa $\int_Z\mathrm d\beta
=\int_{\partial Z}\beta=0$. Wobec tego całka po cyklu
nie rozróżnia reprezentantów jednej klasy de Rhama.
Twierdzenie de Rhama (wynik zewnętrzny) stwierdza,
że całkowanie po cyklach, po utożsamieniu przestrzeni
funkcjonałów na homologii z kohomologią singularną
ze współczynnikami rzeczywistymi, wyznacza naturalny
izomorfizm $H^k_{\rm dR}(M)\simeq H^k(M;\R)$.
Jest to pomost między rachunkiem różniczkowym a topologią.
\end{uwaga}

\begin{twierdzenie}[Dualność Poincarégo — wynik zewnętrzny]
\label{tw:ga10-pd}
Niech $M$ będzie zwartą, spójną i zorientowaną rozmaitością
gładką wymiaru rzeczywistego $m$, bez brzegu.
Dla każdego $0\le k\le m$ parowanie
\[
 H^k(M;\R)\times H^{m-k}(M;\R)\longrightarrow\R,
 \qquad ([\alpha],[\beta])\longmapsto\int_M\alpha\wedge\beta
\]
jest niezdegenerowane. W szczególności $b_k(M)=b_{m-k}(M)$.
\end{twierdzenie}

\begin{proof}[Dowód ostatniej równości z niezdegenerowania]
Formy gładkie reprezentują klasy rzeczywiste przez twierdzenie
de Rhama. Ponieważ obie przestrzenie kohomologii są skończenie
wymiarowe, niezdegenerowane parowanie identyfikuje
$H^k(M;\R)$ z przestrzenią dualną $H^{m-k}(M;\R)^*$.
Przestrzeń i jej dualna mają ten sam wymiar.
Główna część twierdzenia, czyli konstrukcja dualności,
pozostaje wynikiem topologii algebraicznej przyjętym zewnętrznie.
\end{proof}

\begin{definicja}[Klasa dualna do podrozmaitości]
Niech $M$ spełnia założenia twierdzenia~\ref{tw:ga10-pd},
a $Z\subset M$ będzie zwartą, zorientowaną, gładką
podrozmaitością wymiaru $m-2p$.
\emph{Klasą dualną Poincarégo} do $Z$ nazywamy
jedyną klasę $\operatorname{PD}[Z]\in H^{2p}(M;\R)$ taką, że
$\int_Z\eta=\int_M\operatorname{PD}[Z]\wedge\eta$
dla każdej zamkniętej formy $\eta$ stopnia $m-2p$.
Istnienie i jednoznaczność wynikają z dualności Poincarégo.
\end{definicja}

\begin{uwaga}[W przypadku osobliwym]
Cykle algebraiczne mogą być osobliwe, więc nie zawsze są
podrozmaitościami. Dla pododmiany zespolonej czystego wymiaru
$p$ istnieje jednak kanoniczna klasa fundamentalna w homologii
singularnej wymiaru $2p$. Konstrukcja w ogólnym przypadku
wymaga teorii cykli lub homologii Borel–Moore'a; w dalszych
obliczeniach wystarczą gładkie pododmiany i liniowość cykli.
\end{uwaga}

\begin{stwierdzenie}[Intersekcja jako całka]
\label{st:ga10-intersection}
Niech $S$ będzie gładką spójną rzutową powierzchnią nad $\C$,
a $C,D\subset S$ dwiema gładkimi krzywymi przecinającymi się
transwersalnie w skończenie wielu punktach. Wtedy
\[
 C\cdot D=\#(C\cap D)
 =\int_{S^{\rm an}}\operatorname{PD}[C]\wedge\operatorname{PD}[D].
\]
\end{stwierdzenie}
\begin{proof}
W lokalnych współrzędnych holomorficznych wokół każdego
punktu przecięcia krzywe są dwiema transwersalnymi
zespolonymi prostymi. Orientacje zespolone wymuszają dodatni
lokalny znak przecięcia, czyli $+1$. Suma lokalnych znaków
równa się algebraicznej liczbie przecięć, którą dla dwóch
transwersalnych krzywych obliczaliśmy już w
rozdziale~\ref{ch:ga6-powierzchnie}. Z twierdzenia o zgodności
iloczynu kohomologicznego i przecięcia geometrycznego
(zewnętrzny wynik o dualności Poincarégo i klasach przecięć) ta suma jest
równa całce z iloczynu klas dualnych.
\end{proof}

\begin{przyklad}[Dwie proste w $\PP^2$]
Dla dwóch różnych prostych $L_1,L_2\subset\PP^2$ mamy
$L_1\cdot L_2=1$ na mocy Bézouta. Jeśli
$h=\operatorname{PD}[L_1]=\operatorname{PD}[L_2]$,
to $\int_{\PP^2}h^2=1$. Z jednej liczby otrzymujemy
podstawowy rachunek w pierścieniu kohomologii $\PP^2$.
\end{przyklad}

\section{Topologia przestrzeni rzutowej i krzywych}
\label{sec:ga10-projektywne}

\begin{twierdzenie}[Kohomologia przestrzeni rzutowej — wynik standardowy]
\label{tw:ga10-cp-cohom}
Dla $n\ge0$ zachodzi izomorfizm stopniowanych pierścieni
\[
 H^*(\PP^n(\C);\Z)\simeq\Z[h]/(h^{n+1}),
 \qquad \deg h=2,
\]
gdzie $h=c_1(\mathcal O_{\PP^n}(1))$ i
$\int_{\PP^n(\C)}h^n=1$.
\end{twierdzenie}
\begin{proof}[Idea dowodu przez strukturę komórkową]
Przestrzeń $\PP^n(\C)$ otrzymujemy z
$\PP^{n-1}(\C)$, doklejając jedną rzeczywistą komórkę
wymiaru $2n$: chart $Z_n\ne0$ jest zespoloną przestrzenią
afiniczną, a jego dopełnienie to $\PP^{n-1}(\C)$.
Zatem $\PP^n(\C)$ ma po jednej komórce
w wymiarach $0,2,\ldots,2n$ i nie ma komórek
w nieparzystych wymiarach. W komórkowym kompleksie
łańcuchowym wszystkie różniczki są zerowe,
ponieważ przynajmniej jedna z sąsiednich grup jest zerowa.
Stąd $H^{2j}(\PP^n;\Z)\cong\Z$ dla $0\le j\le n$,
a pozostałe grupy dodatnich stopni znikają.
Identyfikacja mnożenia wynika z tego, że hiperpłaszczyzna
reprezentuje $h$, a przecięcie $j$ hiperpłaszczyzn w
położeniu ogólnym jest liniową przestrzenią rzutową
kowymiaru $j$; dla $j=n$ przecięcie jest jednym punktem.
Tym samym $h^j$ generuje $H^{2j}$ i $h^n$ ma całkę równą $1$.
Użyliśmy standardowej zgodności dualności Poincarégo
z iloczynem przecięć z poprzedniej sekcji.
\end{proof}

\begin{stwierdzenie}[Rodzaj krzywej z kohomologii]
Niech $C$ będzie gładką spójną krzywą rzutową nad $\C$
o rodzaju $g$. Wtedy $\chi_{\rm top}(C)=2-2g$ oraz
$\dim_{\C}H^0(C,\Omega^1_C)=g$.
\end{stwierdzenie}
\begin{proof}
Pierwszy wzór wynika z klasyfikacji zwartych zorientowanych
powierzchni (twierdzenie topologiczne): powierzchnia rodzaju $g$
ma rozkład CW z jednym wierzchołkiem, $2g$ krawędziami
i jedną komórką dwuwymiarową. Dlatego jej charakterystyka
Eulera wynosi $1-2g+1=2-2g$.
Drugi wzór otrzymujemy przez identyfikację
$H^0(C,\Omega^1_C)$ ze składnikiem Hodge'a $H^{1,0}$.
Dla krzywej o rodzaju $g$ liczba $b_1$ wynosi $2g$,
a symetria Hodge'a daje $b_1=h^{1,0}+h^{0,1}=2h^{1,0}$.
W szczególności $h^{1,0}=g$, zgodnie z Riemannem--Rochem
z rozdziału~\ref{ch:krzywe-snopy}.
\end{proof}

\begin{center}
\begin{tikzpicture}[>=Stealth,every node/.style={font=\footnotesize}]
 \node[draw=blue!55!black,fill=blue!5,rounded corners=5pt,
       minimum height=18mm,text width=43mm,align=left,inner sep=6pt]
       (curve) at (-3.25,0)
       {Krzywa $C$ rodzaju $g$\\[2pt]
        $b_1=2g$\\[2pt]
        $\chi_{\rm top}=2-2g$};
 \node[draw=green!45!black,fill=green!6,rounded corners=5pt,
       minimum height=18mm,text width=43mm,align=left,inner sep=6pt]
       (forms) at (3.25,0)
       {Formy holomorficzne\\[2pt]
        $h^{1,0}=g$\\[2pt]
        $h^{0,1}=g$};
 \draw[->,thick] (curve.east)--node[above,font=\scriptsize]{Hodge}(forms.west);
\end{tikzpicture}

{\small Rysunek 10.1. Rzeczywista topologia i holomorficzne różniczki
krzywej rzutowej niosą tę samą informację o rodzaju.}\par
\end{center}

\section{Słabe twierdzenie Lefschetza: co zachowuje przekrój hiperpłaszczyzną?}
\label{sec:ga10-weak}

\begin{twierdzenie}[Słabe Lefschetza — wynik zewnętrzny]
\label{tw:ga10-weak}
Niech $X$ będzie gładką, spójną, rzutową odmianą nad $\C$
o wymiarze zespolonym $n\ge2$, a $Y=X\cap H$ gładkim
przekrojem przez hiperpłaszczyznę $H$ w ustalonym zanurzeniu
rzutowym $X\hookrightarrow\PP^N_\C$.
Wówczas dla zanurzenia $i\colon Y\hookrightarrow X$ homomorfizm
\[
 i^*\colon H^k(X^{\rm an};\Z)\longrightarrow
 H^k(Y^{\rm an};\Z)
\]
jest izomorfizmem dla $k<n-1$ i monomorfizmem dla $k=n-1$.
Co więcej, $Y$ jest spójna. Wersja homotopijna mówi, że
$\pi_k(Y)\to\pi_k(X)$ jest izomorfizmem dla $k<n-1$
i surjekcją dla $k=n-1$ (dla odpowiednich punktów bazowych).
\end{twierdzenie}

\begin{uwaga}[Skąd biorą się granice $n-1$?]
Jedna droga do dowodu prowadzi przez teorię Morse'a:
w odpowiednio dobranej funkcji na dopełnieniu hiperpłaszczyzny
pojawiają się uchwyty rzeczywistego wymiaru $n$.
Oznacza to, że dołączenie takich komórek nie zmienia
niskich grup homotopii. Jest to jedynie idea dowodu:
uzasadnienie istnienia właściwej funkcji Morse'a,
a następnie konstrukcja odpowiedniego rozkładu CW,
są zasadniczą treścią twierdzenia zewnętrznego.
\end{uwaga}

\begin{stwierdzenie}[Gładkie powierzchnie w $\PP^3$ są po prostu spójne]
\label{st:ga10-simplyconn}
Jeżeli $S\subset\PP^3_\C$ jest gładką zespoloną
hiperpowierzchnią stopnia $d\ge1$, to
$\pi_1(S^{\rm an})=0$ i $b_1(S)=b_3(S)=0$.
\end{stwierdzenie}
\begin{proof}
Stosujemy wersję homotopijną
Lefschetza~\ref{tw:ga10-weak} do pary
$X=\PP^3_\C$, $Y=S$, dla $n=3$.
Otrzymujemy izomorfizm $\pi_1(S)\simeq\pi_1(\PP^3)$.
Ta ostatnia grupa jest trywialna: $\PP^3$ ma strukturę
CW z jedną komórką w każdym parzystym wymiarze od $0$ do $6$,
a więc jej jednoszkielet jest punktem, co daje trywialność
$\pi_1$. Zatem $H_1(S;\Z)=0$, więc $b_1(S)=0$.
Ponieważ $S$ jest zwartą zorientowaną rozmaitością
rzeczywistą wymiaru cztery, dualność Poincarégo daje
$b_3(S)=b_1(S)=0$.
\end{proof}

\begin{uwaga}[Słabe a twarde Lefschetza]
Słabe twierdzenie porównuje \emph{dwie różne przestrzenie}:
$Y$ i $X$. Twarde twierdzenie Lefschetza z
rozdziału~\ref{ch:ga9-analiza-zespolona} dotyczy natomiast
mnożenia przez klasę Kählerowską \emph{na tej samej przestrzeni}.
Nie wolno z jednego wyniku bezpośrednio wywnioskować drugiego.
\end{uwaga}

\section{Klasy Cherna i topologiczna charakterystyka Eulera}
\label{sec:ga10-chern-euler}

\begin{twierdzenie}[Chern--Gauss--Bonnet — wynik zewnętrzny]
\label{tw:ga10-cgb}
Jeśli $M$ jest zwartą rozmaitością zespoloną wymiaru $n$,
to dla jej holomorficznej wiązki stycznej $T_M$ zachodzi
\[
 \chi_{\rm top}(M)=\int_M c_n(T_M).
\]
Prawa strona oznacza ocenę całkowitej klasy Cherna
$c_n(T_M)\in H^{2n}(M;\Z)$ na klasie fundamentalnej.
\end{twierdzenie}

\begin{uwaga}[Dlaczego ta tożsamość jest zaskakująca?]
Lewa strona jest niezmiennikiem topologicznym,
liczonym np. z rozkładu komórkowego. Prawa strona
pochodzi z lokalnej geometrii wiązki stycznej.
Przejście między nimi jest szczególnym przypadkiem
uogólnionego twierdzenia Gaussa--Bonneta.
Sam rachunek klas Cherna wykonamy już algebraicznie.
\end{uwaga}

\begin{twierdzenie}[Charakterystyka Eulera gładkiej powierzchni stopnia $d$]
\label{tw:ga10-surface-euler}
Niech $S\subset\PP^3_\C$ będzie gładką powierzchnią
określoną jednym jednorodnym równaniem stopnia $d\ge1$.
Wtedy
\[
 \chi_{\rm top}(S)=d(d^2-4d+6),\qquad
 K_S^2=d(d-4)^2.
\]
\end{twierdzenie}
\begin{proof}
Niech $h=c_1(\mathcal O_{\PP^3}(1))$ i oznaczmy przez ten sam
symbol jego ograniczenie do $S$.
Z poprzedniego rozdziału o klasach Cherna znamy ciąg dokładny
wiązki stycznej i normalnej
$0\to T_S\to T_{\PP^3}|_S\to\mathcal O_S(d)\to0$.
Z multiplikatywności całkowitej klasy Cherna i z ciągu Eulera
otrzymujemy
\[
 c(T_S)=\frac{(1+h)^4}{1+dh}
 =1+(4-d)h+(d^2-4d+6)h^2,
\]
przy czym pomijamy składniki stopnia wyższego niż $4$,
ponieważ $S$ ma zespolony wymiar dwa.
Na mocy Chern--Gauss--Bonnet
$\chi_{\rm top}(S)=\int_S c_2(T_S)$.
Z wzoru na stopień hiperpowierzchni
$\int_S h^2=d$ (przecięcie z dwiema ogólnymi
hiperpłaszczyznami) dostajemy pierwszy wzór.

Ponadto przez adjunkcję
$K_S\cong\mathcal O_S(d-4)$, więc jego pierwsza
klasa Cherna jest równa $(d-4)h$ i
$K_S^2=\int_S (d-4)^2h^2=d(d-4)^2$.
\end{proof}

\begin{przyklad}[Pięć stopni i pięć różnych niezmienników]
Dla $d=1,2,3,4,5$ otrzymujemy kolejno
$\chi_{\rm top}(S)=3,4,9,24,55$.
Wartości $K_S^2$ wynoszą odpowiednio $9,8,3,0,5$.
Gdy $d=4$, kanonika jest trywialna, a mimo to
$\chi_{\rm top}(S)=24$; różnica między tymi niezmiennikami
pokazuje, że znikanie $c_1(T_S)$ nie oznacza znikania
$c_2(T_S)$.
\end{przyklad}

\begin{stwierdzenie}[Liczby Bettiego tych powierzchni]
\label{st:ga10-betti}
Przy założeniach twierdzenia~\ref{tw:ga10-surface-euler} mamy
\[
 b_0=b_4=1,\qquad b_1=b_3=0,\qquad
 b_2=d(d^2-4d+6)-2.
\]
\end{stwierdzenie}
\begin{proof}
Spójność daje $b_0=1$, a dualność Poincarégo $b_4=1$.
Stwierdzenie~\ref{st:ga10-simplyconn} daje $b_1=b_3=0$.
Z definicji charakterystyki Eulera
$\chi_{\rm top}(S)=b_0-b_1+b_2-b_3+b_4=2+b_2$,
więc wynik otrzymujemy przez podstawienie
wzoru z twierdzenia~\ref{tw:ga10-surface-euler}.
\end{proof}

\begin{przyklad}[Sześcienna powierzchnia i kwartyk K3]
Gładka powierzchnia kubiczna ma $b_2=7$,
co odpowiada $\PP^2$ rozdmuchanej w sześciu punktach.
Gładka powierzchnia stopnia czwartego ma $b_2=22$.
Dla niej $K_S\cong\mathcal O_S$, a słabe Lefschetza
zapewnia $\pi_1(S)=0$; otrzymujemy więc powierzchnię K3.
Sam niezmiennik $b_2=22$ nie identyfikuje jeszcze całej
struktury zespolonej ani izomorfizmu algebraicznego.
\end{przyklad}

\section{Rozdmuchanie punktu: co dzieje się z topologią?}
\label{sec:ga10-blowup}

\begin{twierdzenie}[Zmiana niezmienników po rozdmuchaniu]
\label{tw:ga10-blowup}
Niech $S$ będzie gładką spójną rzutową powierzchnią nad $\C$,
niech $p\in S(\C)$, a $\pi\colon\widetilde S=\operatorname{Bl}_pS\to S$
rozdmuchaniem punktu z dzielnikiem wyjątkowym $E\simeq\PP^1$.
Wtedy dla kohomologii rzeczywistej zachodzi rozkład
\[
 H^2(\widetilde S;\R)\simeq\pi^*H^2(S;\R)\oplus\R[E],
\]
który jest ortogonalny względem formy przecięć, oraz
$E^2=-1$. Ponadto
\[
 \chi_{\rm top}(\widetilde S)=\chi_{\rm top}(S)+1,
 \quad b_2(\widetilde S)=b_2(S)+1,
 \quad K_{\widetilde S}^2=K_S^2-1.
\]
\end{twierdzenie}
\begin{proof}
Część algebraiczną $E^2=-1$ oraz
$K_{\widetilde S}=\pi^*K_S+E$ udowodniliśmy
w rozdziale~\ref{ch:ga5-rozdmuchania} i
wykorzystaliśmy w rozdziale~\ref{ch:ga6-powierzchnie}.
W szczególności $\pi^*D\cdot E=0$ dla każdego
dywizora $D$ na $S$, a zatem
\[
 K_{\widetilde S}^2=(\pi^*K_S+E)^2=K_S^2-1.
\]

Topologicznie rozdmuchanie nie zmienia przestrzeni poza
punktem $p$. Nad $p$ wstawiamy $E\simeq\PP^1(\C)$,
czyli sferę $S^2$. Korzystamy z addytywności charakterystyki Eulera
z kohomologii o zwartych nośnikach dla skończonych
rozwarstwień algebraicznych oraz z równości tej
charakterystyki ze zwykłą charakterystyką Eulera
dla zespolonych odmian algebraicznych. Otrzymujemy
$\chi_{\rm top}(\widetilde S)
=\chi_{\rm top}(S\setminus\{p\})+\chi_{\rm top}(E)
=\chi_{\rm top}(S)-1+2$.
Ostatnia równość wyjaśnia dokładnie przyrost o jeden.

Dla pełnej struktury grupy $H^2$ stosujemy topologiczny
wzór na rozdmuchanie, który można udowodnić
ciągiem Mayera--Vietorisa dla otoczenia $E$
i jego dopełnienia; wynika z niego
$H^2(\widetilde S;\R)\simeq\pi^*H^2(S;\R)\oplus\R[E]$.
Właściwość ortogonalności i znak nowego składnika
wynikają z $E\cdot\pi^*D=0$ oraz $E^2=-1$;
dla dowolnych klas kohomologicznych są konsekwencją
zgodności formy przecięć z topologicznym rozkładem.
Ponieważ pozostałe liczby Bettiego nie zmieniają się
(co także wynika z tego samego rozkładu Mayera--Vietorisa),
mamy $b_2(\widetilde S)=b_2(S)+1$.
\end{proof}

\begin{uwaga}[Inna postać twierdzenia]
Rozdmuchanie odpowiada różniczkowo-topologicznie sumie
spójnej $S^{\rm an}\#\overline{\PP^2(\C)}$,
gdzie kreska oznacza odwrócenie orientacji.
To standardowy wynik topologii czterowymiarowej,
nie utożsamienie rozdmuchania z algebraiczną sumą dwóch
odmian. Nowy składnik formy przecięć ma znak ujemny,
co jest algebraicznym śladem odwrócenia orientacji.
\end{uwaga}

\begin{przyklad}[Sześć rozdmuchań]
Dla $S=\PP^2$ wiemy, że $b_2(S)=1$ i
$\chi_{\rm top}(S)=3$. Po rozdmuchaniu sześciu
różnych punktów w położeniu ogólnym otrzymujemy
$b_2=7$, $\chi_{\rm top}=9$, $K^2=9-6=3$.
Są to dokładnie liczby gładkiej powierzchni kubicznej.
Uzyskujemy więc topologiczną kontrolę wcześniejszego
rachunku 27 prostych: cała geometria przecięć rozwija się
w siedmiowymiarowej przestrzeni klas dywizorów.
\end{przyklad}

\begin{center}
\begin{tikzpicture}[>=Stealth,every node/.style={font=\footnotesize}]
 \draw[rounded corners,draw=blue!60!black,fill=blue!4] (-4.5,-1.2) rectangle (-1.65,1.0);
 \node at (-3.1,.62) {$S$};
 \draw[fill=blue!65] (-3.2,-.16) circle (1.8pt);
 \node[anchor=west,font=\scriptsize] at (-3.08,-.16) {$p$};
 \draw[->,thick] (-1.42,-.09)--(.02,-.09)
  node[midway,above,font=\scriptsize] {$\operatorname{Bl}_p$};
 \draw[rounded corners,draw=blue!60!black,fill=blue!4] (.3,-1.2) rectangle (4.45,1.0);
 \node at (2.4,.65) {$\widetilde S$};
 \draw[thick,draw=red!70!black] (1.76,-.32) ellipse (.58 and .30);
 \node[anchor=west,font=\scriptsize] at (2.55,-.13) {$E\simeq\PP^1$};
 \node[anchor=west,font=\scriptsize] at (2.55,-.51) {$E^2=-1$};
 \node[font=\scriptsize] at (0,-1.62) {$b_2(\widetilde S)=b_2(S)+1$};
\end{tikzpicture}

{\small Rysunek 10.2. Rozdmuchanie zastępuje punkt krzywą wyjątkową.
Schemat nie przedstawia lokalnego kształtu rzeczywistej czterowymiarowej
rozmaitości; wyróżnia jedynie zmianę danych topologicznych.}\par
\end{center}

\section{Twarde Lefschetza, formy pierwotne i twierdzenie Hodge'a o indeksie}
\label{sec:ga10-hard-index}

W rozdziale~\ref{ch:ga9-analiza-zespolona} poznaliśmy
twarde twierdzenie Lefschetza. Teraz wyprowadzimy z niego
ograniczenia formy przecięć na powierzchni.

\begin{definicja}[Operator Lefschetza i kohomologia pierwotna]
Niech $M$ będzie zwartą rozmaitością Kählerowską
o wymiarze zespolonym $n$ z formą Kählerowską $\omega$.
Oznaczamy $L(\alpha)=[\omega]\smile\alpha$.
Dla $0\le k\le n$ definiujemy
\[
 P^k(M)=\ker\bigl(L^{n-k+1}\colon
 H^k(M;\C)\to H^{2n-k+2}(M;\C)\bigr).
\]
Elementy $P^k(M)$ nazywamy \emph{klasami pierwotnymi}.
Dla powierzchni $S$ i $k=2$ dostajemy po prostu
$P^2(S)=\{\alpha\in H^2(S;\C):[\omega]\smile\alpha=0\}$.
\end{definicja}

\begin{twierdzenie}[Rozkład Lefschetza — konsekwencja twardego Lefschetza]
Niech $M$ będzie jak w poprzedniej definicji.
Dla $0\le k\le n$ zachodzi rozkład
$H^k(M;\C)=\bigoplus_{r\ge0}L^rP^{k-2r}(M)$,
z pominięciem składników o ujemnych wykładnikach.
\end{twierdzenie}
\begin{proof}[Szkic argumentu z jawnym wynikiem zewnętrznym]
Twarde Lefschetza zapewnia izomorfizmy
$L^{n-j}\colon H^j\to H^{2n-j}$.
Z nich, przez rozkład na jądro odpowiedniej potęgi $L$
i obraz operatora $L$ z poprzedniego stopnia,
otrzymuje się sumę $H^k=P^k\oplus LH^{k-2}$.
Indukcja po $k$ daje zapis w tezie. Pełne uzasadnienie
prostopadłości i jednoznaczności tego rozkładu jest
standardowym wnioskiem z teorii reprezentacji $\mathfrak{sl}_2$
operatorów Lefschetza. Ponieważ nie wprowadzaliśmy tej
reprezentacji, tę część traktujemy jako przyjętą
konsekwencję twardego twierdzenia Lefschetza.
\end{proof}

\begin{twierdzenie}[Hodge'a o indeksie — wynik zewnętrzny]
\label{tw:ga10-index}
Niech $S$ będzie gładką spójną rzutową powierzchnią
nad $\C$, niech $H$ będzie dzielnikiem ample, a
$\rho(S)=\dim_{\R}N^1(S)_{\R}$ jej liczbą Picarda.
Wówczas forma przecięć na rzeczywistej przestrzeni klas
dywizorów modulo równoważność numeryczną
$N^1(S)_{\R}$ ma sygnaturę $(1,\rho(S)-1)$.
Równoważnie, jeśli $D\in N^1(S)_{\R}$,
$D\cdot H=0$ i $D\ne0$, to $D^2<0$.
\end{twierdzenie}

\begin{uwaga}[Skąd bierze się ujemny znak?]
Klasa dzielnika ma typ Hodge'a $(1,1)$.
Jeżeli jest prostopadła do dodatniej klasy $[H]$, to jest
klasą pierwotną stopnia dwa. Relacje dwuliniowe
Hodge'a--Riemanna, przyjęte tutaj jako wynik zewnętrzny,
mówią, że na rzeczywistych pierwotnych klasach $(1,1)$
forma $\alpha\mapsto\int_S\alpha\wedge\alpha$ jest ujemnie
określona. Jedyny dodatni kierunek jest reprezentowany
przez $[H]$, gdyż $H^2>0$. Związek tego rozumowania
z klasami algebraicznymi wyjaśnia twierdzenie Lefschetza
o klasach $(1,1)$ z kolejnej sekcji.
\end{uwaga}

\begin{stwierdzenie}[Nierówność Hodge'a dla dywizorów]
\label{st:ga10-hodge-ineq}
Niech $S,H$ spełniają założenia
twierdzenia~\ref{tw:ga10-index}. Dla każdego
$D\in N^1(S)_{\R}$ zachodzi
\[
 (D\cdot H)^2\ge D^2H^2,
\]
a równość występuje dokładnie wtedy, gdy $D$ jest
rzeczywistą wielokrotnością $H$ w $N^1(S)_{\R}$.
\end{stwierdzenie}
\begin{proof}
Połóżmy $a=(D\cdot H)/H^2$ i $D_0=D-aH$.
Z konstrukcji $D_0\cdot H=0$.
Twierdzenie Hodge'a o indeksie daje $D_0^2\le0$,
a równość tylko dla $D_0=0$. Po rozwinięciu
$D_0^2=D^2-(D\cdot H)^2/H^2$ otrzymujemy tezę.
\end{proof}

\begin{przyklad}[Dlaczego forma przecięć nie jest dodatnio określona?]
Na rozdmuchaniu $S=\operatorname{Bl}_p\PP^2$
przestrzeń $N^1(S)_{\R}$ ma bazę $H,E$,
gdzie $H$ to pullback prostej i $E$ jest dzielnikiem
wyjątkowym. Macierz formy przecięć wynosi
$\left(\begin{smallmatrix}1&0\\0&-1\end{smallmatrix}\right)$.
W szczególności $E^2=-1$. Sama efektywność krzywej
nie implikuje więc nieujemnego samoprzecięcia.
Nie wolno jednak w powyższym twierdzeniu podstawiać
$H$ jako pullback prostej: taki $H$ jest nef, ale
nie jest ample, gdyż $H\cdot E=0$.
Aby uzyskać klasę ample, należy wziąć np.
$A=2H-E$, dla której $A^2=3$ oraz $A\cdot E=1$.
\end{przyklad}

\section{Twierdzenie Lefschetza o klasach \texorpdfstring{$(1,1)$}{(1,1)} i problem Hodge'a}
\label{sec:ga10-11}

\begin{definicja}[Grupy Picarda i Nérona--Severiego]
Niech $X$ będzie gładką spójną rzutową odmianą nad $\C$.
Grupa $\operatorname{Pic}(X)$ to grupa izomorfizmów
wiązek liniowych z działaniem iloczynu tensorowego.
Podgrupa $\operatorname{Pic}^0(X)$ składa się z klas
wiązek algebraicznie równoważnych wiązce trywialnej;
grupę ilorazową
$\operatorname{NS}(X)=\operatorname{Pic}(X)/\operatorname{Pic}^0(X)$
nazywamy grupą Nérona--Severiego.
Jej rząd części wolnej to liczba Picarda $\rho(X)$.
\end{definicja}

\begin{twierdzenie}[Lefschetz $(1,1)$ — wynik zewnętrzny]
\label{tw:ga10-11}
Niech $X$ będzie gładką spójną rzutową odmianą nad $\C$.
Wówczas obraz homomorfizmu
$c_1\colon\operatorname{Pic}(X)\to H^2(X^{\rm an};\Z)$
jest dokładnie zbiorem tych klas całkowitych,
których obraz w $H^2(X^{\rm an};\C)$ należy do
$H^{1,1}(X)$, to znaczy
\[
 \operatorname{im}c_1
 =\{\alpha\in H^2(X^{\rm an};\Z):
 \alpha_{\C}\in H^{1,1}(X)\}.
\]
\end{twierdzenie}

\begin{proof}[Wyjaśnienie przez ciąg wykładniczy]
W topologii analitycznej mamy dokładny ciąg snopów
$0\to\Z_X\to\mathcal O_X\xrightarrow{\exp(2\pi i\,\cdot)}
\mathcal O_X^{\times}\to1$.
Pochodzący z niego fragment długiego ciągu kohomologii
ma postać
\[
 H^1(X,\mathcal O_X^{\times})\xrightarrow{c_1}
 H^2(X;\Z)\longrightarrow H^2(X,\mathcal O_X).
\]
Pierwsza grupa jest analityczną grupą Picarda,
a przez GAGA również algebraiczną grupą Picarda.
Z dokładności wynika, że obraz $c_1$ jest jądrem
strzałki do $H^2(X,\mathcal O_X)$.
Na zwartej rozmaitości Kählerowskiej
$H^2(X,\mathcal O_X)\simeq H^{0,2}(X)$ przez Dolbeaulta.
Ponieważ dla całkowitej klasy jej składowe typu $(2,0)$
i $(0,2)$ są ze sobą sprzężone, zerowanie składowej
$(0,2)$ jest równoważne temu, że klasa ma czysty typ $(1,1)$.
To dowodzi tezy po przyjęciu dokładności ciągu wykładniczego,
GAGA i rozkładu Hodge'a, wyjaśniając zarazem użyte wcześniej
w kryterium Kodairy twierdzenie o klasach $(1,1)$.
\end{proof}

\begin{wniosek}[Ograniczenie liczby Picarda]
\label{wn:ga10-rho}
Dla gładkiej spójnej rzutowej powierzchni $S$ nad $\C$
zachodzi $1\le\rho(S)\le h^{1,1}(S)$.
\end{wniosek}
\begin{proof}
Klasa ample, na przykład $c_1(\mathcal O_S(1))$,
jest całkowita i ma niezerowy kwadrat, więc
$\rho(S)\ge1$.
Twierdzenie~\ref{tw:ga10-11} zanurza rzeczywistą
przestrzeń klas dzielników modulo numeryczną równoważność
w rzeczywistą część $H^{1,1}(S)$.
Jej wymiar rzeczywisty jest równy $h^{1,1}(S)$,
co daje nierówność z góry.
\end{proof}

\begin{przyklad}[K3: dwie różne liczby opisujące kohomologię]
Dla gładkiej kwartykowej powierzchni K3 mamy
$b_2=22$. Pokażemy dalej, że $h^{1,1}=20$.
Jednak liczba Picarda może być znacznie mniejsza:
dla bardzo ogólnej gładkiej kwartykowej powierzchni
w $\PP^3$ zachodzi $\rho=1$ (twierdzenie
Noethera--Lefschetza, wynik zewnętrzny).
A zatem istnieje wiele klas typu $(1,1)$ nad $\C$,
które nie są klasami wymiernych kombinacji dywizorów.
To nie przeczy twierdzeniu Lefschetza $(1,1)$:
warunek \emph{całkowitości} jest niezbędny.
\end{przyklad}

\begin{definicja}[Cykl algebraiczny i klasa Hodge'a]
Niech $X$ będzie gładką spójną rzutową odmianą zespoloną.
Dla cyklu algebraicznego kowymiaru $p$ jego obraz
$\operatorname{cl}(Z)\in H^{2p}(X^{\rm an};\Z)$
nazywamy klasą cyklu. \emph{Wymierną klasą Hodge'a}
stopnia $2p$ nazywamy element
$\alpha\in H^{2p}(X;\Q)$, którego obraz zespolony
należy do $H^{p,p}(X)$.
\end{definicja}

\begin{uwaga}[Hipoteza Hodge'a — czego nie możemy twierdzić?]
Każdy cykl algebraiczny kowymiaru $p$ daje klasę
w $H^{p,p}(X)$; wynika to ze zgodności orientacji
zespolonej ze stopniami form. Odwrotne stwierdzenie,
że każda wymierna klasa Hodge'a jest wymierną kombinacją
klas cykli algebraicznych, jest \emph{hipotezą Hodge'a}
i nie jest znane w ogólności. Dla $p=1$ jest twierdzeniem
Lefschetza~\ref{tw:ga10-11}. Nie należy mylić
rozwiązania dla dzielników z rozwiązaniem
w dowolnym kowymiarze.
\end{uwaga}

\section{Diament Hodge'a i pełne rachunki dla hiperpowierzchni}
\label{sec:ga10-hodge-numbers}

Dla zwartej powierzchni Kählerowskiej $S$ zapisujemy
$h^{p,q}(S)=\dim_{\C}H^q(S,\Omega^p_S)$.
Symetria sprzężenia daje $h^{p,q}=h^{q,p}$,
a dualność Serre'a i Poincarégo prowadzi do
$h^{p,q}=h^{2-p,2-q}$. Zwyczajowo przedstawia się te
liczby w \emph{diamencie Hodge'a}.

\begin{stwierdzenie}[Diament Hodge'a gładkiej powierzchni stopnia $d$]
\label{st:ga10-hodge-hypersurf}
Niech $S\subset\PP^3_\C$ będzie gładką powierzchnią
stopnia $d\ge1$. Niech $p_g=h^0(S,K_S)$.
Wtedy $q=h^1(S,\mathcal O_S)=0$ oraz
\[
 p_g=\begin{cases}0,&1\le d\le3,\\
 \binom{d-1}{3},&d\ge4.\end{cases}
\]
Ponadto $h^{1,1}(S)=b_2(S)-2p_g$, gdzie
$b_2(S)$ jest dane wzorem ze
stwierdzenia~\ref{st:ga10-betti}.
\end{stwierdzenie}
\begin{proof}
Z $b_1(S)=0$ i symetrii Hodge'a mamy
$0=b_1=h^{1,0}+h^{0,1}=2h^{0,1}$.
Twierdzenie Dolbeaulta identyfikuje
$h^{0,1}$ z $\dim H^1(S,\mathcal O_S)=q$,
więc $q=0$.

Przez adjunkcję $K_S\simeq\mathcal O_S(d-4)$.
Jeśli $d<4$, wiązka $\mathcal O_S(d-4)$ nie ma
niezerowych przekrojów, ponieważ ma ujemny stopień
na każdej ogólnej sekcji hiperpłaszczyznowej.
Gdy $d\ge4$, korzystamy z ciągu dokładnego
\[
 0\to\mathcal O_{\PP^3}(-4)\longrightarrow
 \mathcal O_{\PP^3}(d-4)\longrightarrow
 \mathcal O_S(d-4)\to0.
\]
Stosujemy standardowe twierdzenie o kohomologii wiązek
$\mathcal O(m)$ na przestrzeni rzutowej (w tym miejscu
przyjęte zewnętrznie w wymiarze trzy), które daje
$H^1(\PP^3,\mathcal O(-4))=0$. Wtedy dokładność
długiego ciągu kohomologii pokazuje, że każdy przekrój
kanoniki jest restrykcją jednorodnego wielomianu
stopnia $d-4$. Przekrojów tych jest
$\binom{(d-4)+3}{3}=\binom{d-1}{3}$.

Wreszcie dla stopnia dwa rozkład Hodge'a ma postać
$b_2=h^{2,0}+h^{1,1}+h^{0,2}$.
Z symetrii Hodge'a oraz identyfikacji
$h^{2,0}=h^0(S,K_S)=p_g$ dostajemy
$h^{1,1}=b_2-2p_g$.
\end{proof}

\begin{przyklad}[Diament K3]
Dla gładkiej kwartykowej powierzchni $S$ w $\PP^3$
mamy $p_g=1$, $q=0$, $b_2=22$ i $h^{1,1}=20$.
Jej diament Hodge'a to
\[
 \begin{array}{ccccc}
 &&1&&\\
 &0&&0&\\
 1&&20&&1\\
 &0&&0&\\
 &&1&&
 \end{array}
\]
(w kolejnych wierszach odczytujemy stopnie $p+q$).
Wyjaśnia to, dlaczego $22$ niezależne klasy w $H^2$
nie są równoważne $22$ niezależnym klasom dywizorów:
klasy typu $(2,0)$ i $(0,2)$ również wnoszą wkład do $b_2$.
\end{przyklad}

\begin{center}
\small
\begin{tabular}{c|r|r|r|r|r}
$d$ & $\chi_{\rm top}$ & $b_2$ & $p_g$ & $h^{1,1}$ & $K_S^2$\\\hline
$1$ & $3$ & $1$ & $0$ & $1$ & $9$\\
$2$ & $4$ & $2$ & $0$ & $2$ & $8$\\
$3$ & $9$ & $7$ & $0$ & $7$ & $3$\\
$4$ & $24$ & $22$ & $1$ & $20$ & $0$\\
$5$ & $55$ & $53$ & $4$ & $45$ & $5$
\end{tabular}

\smallskip
{\small Tabela 10.1. Niezmienniki gładkich powierzchni
stopnia $d$ w $\PP^3$; wszystkie kolumny wynikają
z formuł udowodnionych w tym rozdziale przy wskazanych
twierdzeniach zewnętrznych.}\par
\end{center}

\begin{stwierdzenie}[Kontrola wzoru Noethera]
Dla gładkiej powierzchni $S\subset\PP^3$ stopnia $d$
zachodzi $12\chi(S,\mathcal O_S)=K_S^2+\chi_{\rm top}(S)$.
\end{stwierdzenie}
\begin{proof}
Z obliczonych wyżej liczb Hodge'a wynika
$\chi(S,\mathcal O_S)=h^0-h^1+h^2=1+p_g$;
od strony algebraicznej $h^2(\mathcal O_S)=h^0(K_S)$
z dualności Serre'a. Dla $d\ge4$ dostajemy
$12(1+\binom{d-1}{3})$.
Po podstawieniu wzorów na $K_S^2$ i
$\chi_{\rm top}(S)$ prawa strona jest równa
$d(d-4)^2+d(d^2-4d+6)=2d(d^2-6d+11)$.
Dla $d\ge4$ rozwinięcie dwumianu daje
$12+2(d-1)(d-2)(d-3)=2d(d^2-6d+11)$.
Dla $d=1,2,3$ obie strony wynoszą $12$.
Jest to bezpośrednia kontrola szczególnego przypadku
wzoru Noethera, a nie niezależny dowód pełnego
Hirzebrucha--Riemanna--Rocha.
\end{proof}

\section{Osobliwość w rodzinie i cykl zanikający}
\label{sec:ga10-vanishing}

Topologia gładkiej krzywej może zmienić się, kiedy
w rodzinie pojawi się punkt osobliwy. Jest to ważna
motywacja teorii osobliwości oraz monodromii.

\begin{przyklad}[Lokalne wygładzenie węzła]
\label{prz:ga10-node}
Rozważmy w $\C^2$ rodzinę $C_t=\{(u,v):uv=t\}$,
dla parametrów $t\in\C$ bliskich zeru.
Dla $t=0$ otrzymujemy dwie osie, które przecinają się
transwersalnie w punkcie $(0,0)$. Dla $t\ne0$
równość $v=t/u$ identyfikuje $C_t$ z $\C^*$.
Jeżeli ograniczymy się do małej kuli wokół zera,
zbiór $C_t$ dla odpowiednio małego $t\ne0$ jest
analitycznym pierścieniem (ma typ homotopii okręgu).
Gdy $t\to0$, ten okrąg może skurczyć się do punktu
w osobliwym włóknie. Nazywamy go \emph{cyklem zanikającym}.
\end{przyklad}

\begin{proof}[Dlaczego pojawia się pierścień?]
Wybierzmy promień $\varepsilon>0$ i ograniczmy $C_t$
do kuli $B_\varepsilon=\{|u|^2+|v|^2<\varepsilon^2\}$.
Dla $t\ne0$ podstawienie $v=t/u$ prowadzi do
$|u|^2+|t|^2/|u|^2<\varepsilon^2$.
Przy $0<2|t|<\varepsilon^2$ nierówność ma postać
$r_-<|u|<r_+$ dla pewnych $0<r_-<r_+$.
Istotnie, po podstawieniu $s=|u|^2$ otrzymujemy
nierówność kwadratową $s^2-\varepsilon^2s+|t|^2<0$
o dwóch dodatnich pierwiastkach.
Jest to obszar pierścieniowy w płaszczyźnie $u$,
a więc odkształca się do okręgu $|u|=\sqrt{|t|}$.
Ten okrąg jest cyklem zanikającym.
\end{proof}

\begin{center}
\begin{tikzpicture}[>=Stealth,every node/.style={font=\footnotesize}]
 \node at (-2.95,1.48) {$uv=0$};
 \draw[->,gray] (-4.25,0)--(-1.55,0);
 \draw[->,gray] (-2.95,-1.18)--(-2.95,1.18);
 \draw[very thick,blue!75!black] (-4.05,0)--(-1.85,0);
 \draw[very thick,blue!75!black] (-2.95,-1.03)--(-2.95,1.04);
 \fill[red!70!black] (-2.95,0) circle (2.5pt);
 \draw[->,thick] (-1.16,.06)--(.98,.06)
  node[midway,above,font=\scriptsize] {wygładzenie};
 \node at (2.72,1.48) {$uv=t\ne0$};
 \draw[blue!75!black,thick] (2.72,0) ellipse (1.18 and .67);
 \draw[blue!75!black,dashed] (2.72,0) ellipse (.49 and .31);
 \draw[->,red!70!black,thick] (2.77,-.29) arc (-80:135:.42 and .30);
 \node[font=\scriptsize,align=center] at (2.72,-1.17)
  {cykl zanikający\\(typ homotopii $S^1$)};
\end{tikzpicture}

{\small Rysunek 10.3. Schemat topologiczny węzła i jego wygładzenia.
Prawy rysunek jest modelem pierścienia, a nie rzeczywistym wykresem
pełnej krzywej zespolonej $uv=t$ w czterech wymiarach rzeczywistych.}\par
\end{center}

\begin{przyklad}[Rzutowa krzywa węzłowa i normalizacja]
Rozważmy krzywą kubiczną
$C_0=\{Y^2Z=X^2(X+Z)\}\subset\PP^2_\C$.
W punkcie $[0:0:1]$ ma ona węzeł: w charcie $Z=1$
równanie brzmi $y^2=x^2(x+1)$, a jego częścią
jednorodną najniższego stopnia jest $y^2-x^2$,
czyli iloczyn dwóch różnych form liniowych.
Krzywa jest nierozkładalna: w charcie afinicznym wielomian
$y^2-x^2(x+1)$ nie rozkłada się w $\C[x,y]$, ponieważ
$x+1$ nie jest kwadratem w $\C[x]$, a jako wielomian
moniczny stopnia dwa w $y$ mógłby rozłożyć się tylko
po istnieniu pierwiastka w $\C[x]$.
Jej normalizacja ma jednak rodzaj zero. Istotnie, na otwartym podzbiorze
$x\ne0$ możemy wprowadzić $s=y/x$; wtedy
$x=s^2-1$ oraz $y=s(s^2-1)$.
Jednorodnie zapisujemy tę parametryzację jako morfizm
$\nu\colon\PP^1\to C_0$ dany wzorem
$\nu([S:T])=[T(S^2-T^2):S(S^2-T^2):T^3]$.
Wzór ma sens także dla $T=0$, nie ma wspólnych zer wszystkich
trzech współrzędnych i spełnia równanie kubiki.
Jest to morfizm biwymierny z gładkiej krzywej $\PP^1$;
ponieważ jest skończony (jest właściwy i ma skończone włókna),
stanowi normalizację $C_0$ na mocy własności uniwersalnej
normalizacji krzywych. Użyliśmy tu standardowych faktów
o skończoności właściwego morfizmu quasi-skończonego.
Morfizm rozdziela dwie gałęzie węzła:
$s=1$ i $s=-1$ mają ten sam obraz $[0:0:1]$.
Zatem topologicznie $C_0$ otrzymujemy z $S^2$
przez sklejenie dwóch różnych punktów, a więc
$\chi_{\rm top}(C_0)=2-1=1$.

Dla małego ogólnego $t\ne0$ równanie
$Y^2Z=X^2(X+Z)+tZ^3$ definiuje gładką
krzywą kubiczną, czyli krzywą rodzaju jeden.
Jej charakterystyka Eulera jest równa zero.
Zmiana $\chi_{\rm top}$ z $1$ na $0$ jest
obrazem lokalnej zmiany w węźle: w miejsce
skrzyżowania dwóch gałęzi pojawia się pierścień.
\end{przyklad}

\begin{uwaga}[Monodromia i wzór Picarda--Lefschetza]
Przechodząc z parametrem $t$ raz wokół zera,
możemy przenosić cykle gładkich włókien.
Otrzymany automorfizm homologii jest \emph{monodromią}.
Dla pojedynczego zwyczajnego punktu podwójnego
w odpowiednim wymiarze opisuje go wzór Picarda--Lefschetza:
transformacja różni się od identyczności o składnik
wyznaczony przez przecięcie z cyklem zanikającym
(z dokładnym znakiem zależnym od wymiaru i przyjętych
konwencji orientacji). Pełna konstrukcja monodromii
wymaga teorii fibracji Milnora i pozostaje poza
zakresem bieżącego rozdziału. Lokalny rachunek $uv=t$
pokazuje jednak geometryczną ideę stojącą za tym wzorem.
\end{uwaga}

\section{Jakie ograniczenia topologiczne są szczególnie użyteczne?}
\label{sec:ga10-konsekwencje}

\begin{stwierdzenie}[Trzy testy dla gładkiej rozmaitości rzutowej]
Niech $X$ będzie gładką spójną rzutową odmianą nad $\C$
o wymiarze zespolonym $n$. Wtedy:
\begin{enumerate}[label=\arabic*)]
\item jej nieparzyste liczby Bettiego są parzyste;
\item dla każdego $k<n$ zachodzi $b_k(X)\le b_{k+2}(X)$;
\item jeżeli $n=2$, to forma przecięć na $N^1(X)_{\R}$
ma dokładnie jeden dodatni kierunek, a pozostałe
$\rho(X)-1$ kierunków są ujemne.
\end{enumerate}
\end{stwierdzenie}
\begin{proof}
Punkt pierwszy wyprowadziliśmy z symetrii Hodge'a
we wniosku~\ref{wn:ga9-odd-betti}. Drugi wynika z
injektywności mnożenia przez klasę Kählerowską
w stopniach $k<n$, co zostało uzasadnione z twardego
Lefschetza w rozdziale~\ref{ch:ga9-analiza-zespolona}.
Trzeci jest twierdzeniem Hodge'a o indeksie
\ref{tw:ga10-index}. Każde założenie rzutowości
jest tu istotne co najmniej jako gwarancja struktury
Kählerowskiej oraz odpowiednich własności klas dzielników.
\end{proof}

\begin{przyklad}[Powierzchnia Hopfa jako kontrprzykład]
Powierzchnia Hopfa z rozdziału~\ref{ch:ga9-analiza-zespolona}
ma $b_1=1$. Nie może więc być Kählerowska,
a w szczególności nie może być gładką rzutową
powierzchnią nad $\C$. Jest to ograniczenie
czysto topologiczne, wykrywające brak zanurzenia
rzutowego bez rozwiązywania równań wielomianowych.
\end{przyklad}

\begin{uwaga}[Ostrożność: nie każdy niezmiennik wyznacza odmianę]
Dwie powierzchnie o jednakowych liczbach Bettiego,
charakterystyce Eulera i liczbach Cherna nie muszą być
izomorficzne algebraicznie ani biholomorficzne.
Nawet pełna topologia różniczkowa nie determinuje
w ogólności struktury zespolonej. Dla K3 istnieją
wielowymiarowe rodziny nieizomorficznych struktur,
a zmieniająca się część geometrii jest widoczna
w okresach form holomorficznych i w położeniu krat
klas algebraicznych wewnątrz $H^2$.
\end{uwaga}

\section{Synteza i mapa zależności między twierdzeniami}
\label{sec:ga10-koniec}

W tym rozdziale powiązaliśmy trzy poziomy opisu.
Równania wielomianowe dostarczają wiązek, klas Cherna
oraz cykli; analiza Kählerowska dostarcza form
harmonicznych, rozkładu Hodge'a i operatorów Lefschetza;
topologia dostarcza klas fundamentalnych, kohomologii
singularnej, charakterystyki Eulera i form przecięć.

\begin{center}
\begin{tikzpicture}[>=Stealth,every node/.style={font=\footnotesize,align=center},
  bx/.style={draw,rounded corners=3pt,text width=40mm,minimum height=14mm,inner sep=4pt}]
 \node[bx,fill=blue!4] (alg) at (-3.6,.9)
   {geometria algebraiczna\\cykle, dywizory, klasy Cherna};
 \node[bx,fill=green!4] (an) at (3.5,.9)
   {analiza zespolona\\Dolbeault, Hodge, Kähler};
 \node[bx,fill=orange!5] (top) at (0,-1.45)
   {topologia\\kohomologia, przecięcia, Euler};
 \draw[<->,thick] (alg)--node[above,font=\scriptsize] {GAGA, $(1,1)$} (an);
 \draw[->,thick] (alg)--node[below left,font=\scriptsize] {klasa cyklu} (top);
 \draw[->,thick] (an)--node[below right,font=\scriptsize] {de Rham, Lefschetz} (top);
\end{tikzpicture}

{\small Rysunek 10.4. Powiązania trzech języków geometrii.
Strzałki oznaczają określone konstrukcje i twierdzenia
porównawcze, a nie równość odpowiednich kategorii.}\par
\end{center}

\begin{uwaga}[Granice samowystarczalności]
Dowiedliśmy szczegółowo formuł dla $\chi_{\rm top}$,
$K^2$, liczb Bettiego i części liczb Hodge'a
gładkich powierzchni stopnia $d$ w $\PP^3$,
pokazaliśmy zmianę tych niezmienników przy rozdmuchaniu,
wyprowadziliśmy nierówność Hodge'a dla dywizorów,
wyjaśniliśmy ciąg wykładniczy oraz obliczyliśmy
homotopię lokalnego wygładzenia węzła.
Jako twierdzenia zewnętrzne pozostają: dualność Poincarégo,
twierdzenie de Rhama, słabe i twarde Lefschetza,
twierdzenie Hodge'a o indeksie i relacje Hodge'a--Riemanna,
Chern--Gauss--Bonnet, GAGA, Dolbeault i ogólne twierdzenia
porównawcze o klasach cykli. Zostały one użyte z jawnymi
założeniami, a nie ukryte w rachunkowych przeskokach.
\end{uwaga}

% Podstawowe pozycje do dalszej lektury i weryfikacji:
% W. Fulton, Intersection Theory.
% C. Voisin, Hodge Theory and Complex Algebraic Geometry I.
% P. Griffiths, J. Harris, Principles of Algebraic Geometry.
% R. Hartshorne, Algebraic Geometry.
% A. Hatcher, Algebraic Topology (duality, cohomology, Euler).
% Stacks Project: https://stacks.math.columbia.edu/

